% =====================================================================
% pertemuan-06.tex
% Pertemuan 6 - Context-Free Grammar: BNF/EBNF, Derivasi, Pohon Sintaks
% Sub-CPMK 3.1, 3.4 - Praktik terbimbing
% =====================================================================
\documentclass[aspectratio=169,10pt,t]{beamer}
\usepackage{config/beamerthemeubb}

\UBBmeeting{6}

\begin{document}

\UBBtitleframe{6}{Context-Free Grammar: BNF/EBNF, Derivasi, dan Pohon Sintaks}{Praktik terbimbing}

% ---------------------------------------------------------------------
\begin{frame}{Capaian Pertemuan}
  \textbf{CPMK induk:} CPMK 03 --- Merancang tata bahasa
  (\textit{grammar}) bebas konteks untuk bahasa sederhana.

  \vspace{3pt}
  \begin{block}{Sub-CPMK yang dinilai pada pertemuan ini}
    \footnotesize
    \begin{itemize}
      \setlength{\itemsep}{2pt}
      \item \textbf{Sub-CPMK 3.1:} Membuat berkas grammar BNF dan EBNF
            untuk subset bahasa C, serta menguji derivasi dan pohon
            sintaks yang dihasilkan.
      \item \textbf{Sub-CPMK 3.4:} Membuat berkas grammar bebas konteks
            untuk subset bahasa C yang mencakup deklarasi, ekspresi,
            dan pernyataan, sebagai dasar parser dan proyek akhir.
    \end{itemize}
  \end{block}
\end{frame}

% ---------------------------------------------------------------------
\begin{frame}{Agenda Pertemuan (150 menit)}
  \small
  \begin{tabular}{@{}p{1.4cm}p{12.3cm}@{}}
    \toprule
    \textbf{Menit} & \textbf{Kegiatan} \\
    \midrule
    30 & Ceramah: CFG, BNF, dan bahasa yang dibangkitkan \\
    30 & Latihan: konversi EBNF ke BNF \\
    40 & Praktik: menyusun grammar deklarasi, ekspresi, pernyataan \\
    30 & Praktik: lima derivasi dan dua pohon urai \\
    20 & Peer review berkas grammar dan penutup \\
    \bottomrule
  \end{tabular}

  \vspace{6pt}
  {\scriptsize\color{ubbblue!75}Porsi praktik 120 menit. Keluaran
  pertemuan adalah berkas grammar dan bukti derivasi/pohon, bukan kode
  parser.}
\end{frame}

% ---------------------------------------------------------------------
\begin{frame}{Posisi Pertemuan pada Rantai Proyek}
  \UBBchain{2}

  \vspace{2pt}
  \begin{itemize}
    \item Pertemuan ini membangun simpul \textbf{grammar}: aturan yang
          menentukan bentuk program Subset-C.
    \item Grammar menjadi masukan langsung untuk analisis sintaksis,
          pembersihan rekursi kiri, dan parser pada Pertemuan 7--9.
    \item Berkas EBNF/BNF adalah fondasi proyek akhir.
  \end{itemize}
\end{frame}

% ---------------------------------------------------------------------
\begin{frame}{CFG dan Notasi BNF}
  \begin{block}{\textit{Context-Free Grammar} (CFG)}
    \small
    \(G=(V,T,P,S)\), dengan \(V\) himpunan nonterminal, \(T\) himpunan
    terminal, \(P\) aturan \(A \to \alpha\) (\(A \in V\),
    \(\alpha \in (V \cup T)^{*}\)), dan \(S \in V\) simbol awal.
  \end{block}

  \begin{itemize}
    \item Aturan mengganti satu nonterminal tanpa memedulikan konteks.
    \item BNF menulis alternatif dengan \(\mid\):
          \(A \to \alpha_1 \mid \alpha_2\).
    \item Nonterminal ditulis miring; terminal dalam tanda kutip atau
          huruf mesin.
    \item Produksi \(\varepsilon\) ditulis eksplisit.
  \end{itemize}
\end{frame}

% ---------------------------------------------------------------------
\begin{frame}{Contoh CFG Aritmetika}
  \[
    E \to E + T \mid T, \quad
    T \to T * F \mid F, \quad
    F \to (E) \mid \code{id}.
  \]

  \vspace{4pt}
  \begin{itemize}
    \item \(E\), \(T\), dan \(F\) menyatakan ekspresi, term, dan faktor.
    \item Terminal: \code{id}, \code{+}, \code{*}, \code{(}, \code{)}.
    \item Bahasa yang dibangkitkan:
          \(L(G)=\{w \in T^{*} \mid S \Rightarrow^{*} w\}\).
    \item Contoh lain: \(S \to aSb \mid \varepsilon\) membangkitkan
          \(\{a^n b^n \mid n \ge 0\}\).
  \end{itemize}

  \vspace{4pt}
  {\footnotesize\color{ubbblue!80}CFG tidak memeriksa tipe atau
  deklarasi; pemeriksaan makna dilakukan pada fase semantik.}
\end{frame}

% ---------------------------------------------------------------------
\begin{frame}{Mengapa CFG Cocok untuk Sintaksis}
  \begin{enumerate}
    \setlength{\itemsep}{4pt}
    \item Dapat menyatakan struktur bersarang: kurung, blok, dan
          pemanggilan fungsi.
    \item Dapat diurai secara efisien dengan parser LL atau LR.
    \item Memiliki teori matang untuk analisis dan transformasi.
  \end{enumerate}

  \vspace{6pt}
  \begin{block}{Batas CFG}
    \footnotesize
    CFG memastikan bentuk program benar, tetapi tidak dapat memeriksa
    apakah variabel sudah dideklarasikan atau apakah tipe operand
    cocok. Pemeriksaan itu memerlukan tabel simbol dan aturan semantik.
  \end{block}
\end{frame}

% ---------------------------------------------------------------------
\begin{frame}{Notasi EBNF}
  \small
  \setlength{\tabcolsep}{4pt}
  \begin{tabular}{@{}p{2.6cm}p{3.4cm}p{7.6cm}@{}}
    \toprule
    \textbf{Notasi} & \textbf{Arti} & \textbf{Contoh} \\
    \midrule
    \(X \mid Y\) & pilihan & \texttt{int | float} \\
    \([X]\) & opsional & \texttt{[else stmt]} \\
    \(\{X\}\) & nol atau lebih & \texttt{\{stmt\}} \\
    \((X)\) & pengelompokan & \texttt{(a | b) c} \\
    \bottomrule
  \end{tabular}

  \vspace{5pt}
  \begin{itemize}
    \item EBNF lebih ringkas dan mudah dibaca.
    \item Parser generator seperti Bison memakai BNF.
    \item EBNF dikonversi ke BNF sebelum implementasi.
  \end{itemize}
\end{frame}

% ---------------------------------------------------------------------
\begin{frame}{Konversi EBNF ke BNF}
  \begin{enumerate}
    \setlength{\itemsep}{1pt}
    \item Pilihan \(X \mid Y\) menjadi dua aturan terpisah.
    \item Opsional \([X]\) menjadi \(X\) atau \(\varepsilon\).
    \item Pengulangan \(\{X\}\) menjadi nonterminal bantu rekursif.
    \item Pengelompokan \((X)\) menjadi nonterminal baru bila perlu.
  \end{enumerate}

  \vspace{3pt}
  \begin{block}{Contoh}
    \footnotesize
    \texttt{params = param \{"," param\}} menjadi\\
    \(\text{params} \to \text{param}\ \text{params\_tail}\),\quad
    \(\text{params\_tail} \to \texttt{"," param}\ \text{params\_tail}
     \mid \varepsilon\).\\[2pt]
    \texttt{returnstmt = "return" [expr] ";"} menjadi\\
    \(\text{returnstmt} \to \texttt{"return"}\ \text{opt\_expr}\
     \texttt{";"}\),\quad
    \(\text{opt\_expr} \to \text{expr} \mid \varepsilon\).
  \end{block}
\end{frame}

% ---------------------------------------------------------------------
\begin{frame}[fragile]{Berkas \file{grammar.ebnf}}
  \begin{columns}[T,onlytextwidth]
    \begin{column}{0.58\textwidth}
      \lstinputlisting[language=EBNF,firstline=1,lastline=14,
        firstnumber=1,
        caption={Bagian awal \file{grammar.ebnf}},
        label={lst:p6-ebnf}]{../buku_ajar/code/subsetc/spec/grammar.ebnf}
    \end{column}
    \begin{column}{0.38\textwidth}
      \begin{itemize}
        \item \(\{\}\) untuk daftar, \([]\) untuk bagian opsional.
        \item Satu aturan per konstruksi.
        \item Nama nonterminal mencerminkan peran.
        \item Notasi ringkas untuk dibaca manusia.
      \end{itemize}
      \vspace{2pt}
      {\scriptsize\color{ubbblue!75}Berkas ini harus dikonversi sebelum
      dipakai parser generator.}
    \end{column}
  \end{columns}
\end{frame}

% ---------------------------------------------------------------------
\begin{frame}[fragile]{Berkas \file{grammar.bnf}}
  \begin{columns}[T,onlytextwidth]
    \begin{column}{0.60\textwidth}
      \lstinputlisting[language=EBNF,firstline=5,lastline=18,
        firstnumber=5,
        caption={Hasil konversi pada \file{grammar.bnf}},
        label={lst:p6-bnf}]{../buku_ajar/code/subsetc/spec/grammar.bnf}
    \end{column}
    \begin{column}{0.36\textwidth}
      \begin{itemize}
        \item Simbol bantu berakhiran \texttt{\_tail}.
        \item Bagian opsional memakai awalan \texttt{opt\_}.
        \item Pengulangan menjadi nonterminal rekursif.
        \item Konvensi nama memudahkan penelusuran asal aturan.
      \end{itemize}
    \end{column}
  \end{columns}
\end{frame}

% ---------------------------------------------------------------------
\begin{frame}{Derivasi Leftmost}
  \small
  Untuk grammar aritmetika dan string
  \(\code{id}+\code{id}*\code{id}\):
  \[
  \begin{aligned}
    E &\Rightarrow E+T
      \Rightarrow T+T
      \Rightarrow F+T
      \Rightarrow \code{id}+T\\
    &\Rightarrow \code{id}+T*F
      \Rightarrow \code{id}+F*F
      \Rightarrow \code{id}+\code{id}*F\\
    &\Rightarrow \code{id}+\code{id}*\code{id}.
  \end{aligned}
  \]

  \vspace{4pt}
  \begin{itemize}
    \item Leftmost selalu mengganti nonterminal paling kiri.
    \item Rightmost mengganti nonterminal paling kanan.
    \item Pada grammar tidak ambigu, keduanya menghasilkan pohon yang
          sama.
  \end{itemize}
\end{frame}

% ---------------------------------------------------------------------
\begin{frame}{Bentuk Kalimat dan Pohon Urai}
  \begin{block}{Bentuk kalimat (\eng{sentential form})}
    \small
    Setiap string yang muncul selama derivasi. Bentuk kalimat dapat
    memuat campuran terminal dan nonterminal; kalimat
    (\eng{sentence}) hanya memuat terminal.
  \end{block}

  \begin{itemize}
    \item Pada derivasi sebelumnya, \(E+T*F\) adalah bentuk kalimat.
    \item \(\code{id}+\code{id}*\code{id}\) adalah kalimat.
    \item Akar pohon adalah simbol awal; anak simpul mengikuti ruas
          kanan aturan.
    \item Daun dibaca kiri ke kanan membentuk string turunan.
  \end{itemize}
\end{frame}

% ---------------------------------------------------------------------
\begin{frame}{Pohon Urai untuk \texttt{id + id * id}}
  \begin{center}
    \begin{tikzpicture}[
        nd/.style={draw=ubbblue!40,fill=ubbblue!7,rounded corners=2pt,
          minimum width=0.52cm,minimum height=0.46cm,font=\scriptsize,
          inner sep=1pt},
        ar/.style={draw=ubbblue!55,thick}
      ]
      \node[nd] (E) at (0,0) {\(E\)};
      \node[nd] (E1) at (-2.2,-0.8) {\(E\)};
      \node[nd] (p) at (0,-0.8) {\(+\)};
      \node[nd] (T1) at (2.0,-0.8) {\(T\)};
      \node[nd] (T2) at (-2.2,-1.6) {\(T\)};
      \node[nd] (F1) at (-2.2,-2.4) {\(F\)};
      \node[nd] (id1) at (-2.2,-3.2) {\code{id}};
      \node[nd] (T3) at (1.2,-1.6) {\(T\)};
      \node[nd] (s) at (2.4,-1.6) {\(*\)};
      \node[nd] (F2) at (3.6,-1.6) {\(F\)};
      \node[nd] (F3) at (1.2,-2.4) {\(F\)};
      \node[nd] (id2) at (1.2,-3.2) {\code{id}};
      \node[nd] (F4) at (3.6,-2.4) {\(F\)};
      \node[nd] (id3) at (3.6,-3.2) {\code{id}};
      \foreach \a/\b in {E/E1,E/p,E/T1,T2/F1,F1/id1,
        T1/T3,T1/s,T1/F2,T3/F3,F3/id2,F2/F4,F4/id3}{
        \draw[ar] (\a) -- (\b);
      }
      \draw[ar] (E1) -- (T2);
    \end{tikzpicture}
  \end{center}

  \vspace{1pt}
  {\footnotesize Struktur pohon menunjukkan bahwa perkalian lebih dalam
  daripada penjumlahan, sesuai prioritas operator.}
\end{frame}

% ---------------------------------------------------------------------
\begin{frame}{Pohon Urai dibandingkan AST}
  \small
  \setlength{\tabcolsep}{4pt}
  \begin{tabular}{@{}p{3.2cm}p{5.2cm}p{5.0cm}@{}}
    \toprule
    \textbf{Aspek} & \textbf{Pohon urai} & \textbf{AST} \\
    \midrule
    Isi & Semua simbol tata bahasa & Konstruksi bermakna \\
    Ukuran & Lebih besar & Lebih ringkas \\
    Detail & Kurung, titik koma, kata kunci & Hanya operand dan operator \\
    Penggunaan & Analisis sintaksis & Analisis semantik dan kode \\
    \bottomrule
  \end{tabular}

  \vspace{5pt}
  {\footnotesize\color{ubbblue!80}AST menghilangkan detail sintaksis
  yang tidak diperlukan fase berikutnya. Konversi dapat dilakukan saat
  parser melakukan \eng{reduce}.}
\end{frame}

% ---------------------------------------------------------------------
\begin{frame}{AST untuk \texttt{a + b * c}}
  \begin{columns}[T,onlytextwidth]
    \begin{column}{0.46\textwidth}
      \centering
      \begin{tikzpicture}[
          nd/.style={draw=ubbblue!40,fill=ubbblue!7,rounded corners=2pt,
            minimum width=0.95cm,minimum height=0.52cm,font=\scriptsize},
          ar/.style={-{Stealth[length=1.4mm]},draw=ubbblue!55,thick}
        ]
        \node[nd] (plus) at (0,0) {\(+\)};
        \node[nd] (a) at (-1.3,-0.9) {\code{a}};
        \node[nd] (mul) at (1.2,-0.9) {\(*\)};
        \node[nd] (b) at (0.5,-1.8) {\code{b}};
        \node[nd] (c) at (1.9,-1.8) {\code{c}};
        \draw[ar] (plus) -- (a);
        \draw[ar] (plus) -- (mul);
        \draw[ar] (mul) -- (b);
        \draw[ar] (mul) -- (c);
      \end{tikzpicture}
      \par {\scriptsize AST ekspresi}
    \end{column}
    \begin{column}{0.50\textwidth}
      \begin{block}{Notasi prefix}
        \centering
        \texttt{(+ a (* b c))}
      \end{block}
      \begin{itemize}
        \item Tidak ada tanda kurung sintaksis.
        \item Struktur operator langsung terlihat.
        \item Akar \(+\) menunjukkan penjumlahan sebagai operasi
              terluar.
      \end{itemize}
    \end{column}
  \end{columns}
\end{frame}

% ---------------------------------------------------------------------
\begin{frame}{Hierarki Nonterminal Grammar Subset-C}
  \small
  \setlength{\tabcolsep}{4pt}
  \begin{tabular}{@{}p{3.4cm}p{9.8cm}@{}}
    \toprule
    \textbf{Nonterminal} & \textbf{Tanggung jawab} \\
    \midrule
    \texttt{program} & Daftar deklarasi tingkat atas \\
    \texttt{topdecl} & Deklarasi fungsi atau prototipe \\
    \texttt{stmt} & Deklarasi lokal, blok, dan pernyataan kontrol \\
    \texttt{decl} & Tipe, identifier, dan inisialisasi opsional \\
    \texttt{expr} & Ekspresi; titik masuk ke tingkat prioritas \\
    \texttt{assign} & Penugasan atau ekspresi logika \\
    \texttt{logor}--\texttt{term} & Satu tingkat untuk satu kelompok operator \\
    \texttt{primary} & Literal, identifier, panggilan, kurung \\
    \bottomrule
  \end{tabular}

  \vspace{4pt}
  {\footnotesize\color{ubbblue!80}Satu nonterminal untuk satu tingkat
  prioritas; operator berprioritas lebih tinggi berada lebih dalam
  pohon.}
\end{frame}

% ---------------------------------------------------------------------
\begin{frame}{Aturan Deklarasi dan Fungsi}
  \small
  \[
  \begin{aligned}
    \text{topdecl} &\to \text{type}\ \code{ID}\ \texttt{"("}\
      \text{opt\_params}\ \texttt{")"}\ \text{topdecl\_tail},\\
    \text{topdecl\_tail} &\to \text{block} \mid \texttt{";"},\\
    \text{opt\_params} &\to \text{params} \mid \varepsilon,\\
    \text{params} &\to \text{param}\ \text{params\_tail},\\
    \text{params\_tail} &\to \texttt{"," param}\ \text{params\_tail}
      \mid \varepsilon,\\
    \text{param} &\to \text{type}\ \code{ID}.
  \end{aligned}
  \]

  \vspace{3pt}
  \begin{itemize}
    \item Fungsi dengan tubuh blok atau prototipe tanpa tubuh.
    \item Parameter opsional ditangani \texttt{opt\_params}.
    \item Daftar parameter memakai rekursi kanan.
  \end{itemize}
\end{frame}

% ---------------------------------------------------------------------
\begin{frame}{Aturan Pernyataan}
  \small
  \[
  \begin{aligned}
    \text{stmt} &\to \text{decl}\ \texttt{";"} \mid \text{block}
      \mid \text{ifstmt} \mid \text{whilestmt}\\
    &\quad \mid \text{forstmt} \mid \text{returnstmt}
      \mid \text{breakstmt} \mid \text{continuestmt}\\
    &\quad \mid \text{expr}\ \texttt{";"} \mid \text{printstmt}
      \mid \texttt{";"}.
  \end{aligned}
  \]

  \vspace{4pt}
  \begin{itemize}
    \item Pemisahan ini memudahkan parser memilih cabang berdasarkan
          token awal.
    \item \texttt{if} langsung memilih \texttt{ifstmt}; \texttt{while}
          memilih \texttt{whilestmt}.
    \item Deklarasi lokal dan blok tetap berada dalam \texttt{stmt}.
  \end{itemize}
\end{frame}

% ---------------------------------------------------------------------
\begin{frame}{Aturan Ekspresi Berprioritas}
  \small
  \[
  \begin{aligned}
    \text{expr} &\to \text{assign},\\
    \text{assign} &\to \code{ID}\ \texttt{"="}\ \text{assign}
      \mid \text{logor},\\
    \text{logor} &\to \text{logand}\ \text{logor\_tail},\\
    \text{logand} &\to \text{equality}\ \text{logand\_tail},\\
    \text{equality} &\to \text{relation}\ \text{equality\_tail},\\
    \text{relation} &\to \text{additive}\ \text{relation\_tail},\\
    \text{additive} &\to \text{term}\ \text{additive\_tail},\\
    \text{term} &\to \text{unary}\ \text{term\_tail},\\
    \text{unary} &\to \texttt{"!"}\ \text{unary}
      \mid \texttt{"-"}\ \text{unary}
      \mid \texttt{"+"}\ \text{unary} \mid \text{primary}.
  \end{aligned}
  \]

  \vspace{1pt}
  {\scriptsize\color{ubbblue!75}Prioritas naik dari \texttt{logor} ke
  \texttt{term}; unary dan primary berada paling dalam.}
\end{frame}

% ---------------------------------------------------------------------
\begin{frame}[fragile]{Contoh: Prioritas dan Tanda Kurung pada \file{arith.sc}}
  \begin{columns}[T,onlytextwidth]
    \begin{column}{0.46\textwidth}
      \lstinputlisting[language=C,firstline=1,lastline=7,
        firstnumber=1,
        caption={Program uji prioritas \file{arith.sc}},
        label={lst:p6-arith}]{../buku_ajar/code/subsetc/examples/arith.sc}
    \end{column}
    \begin{column}{0.50\textwidth}
      \begin{itemize}
        \item \texttt{2 + 3 * 4}: akar \(+\) dengan anak kanan \(*\).
        \item \texttt{(2 + 3) * 4}: akar \(*\) dengan anak kiri \(+\).
        \item Perbedaan berasal dari tanda kurung dan tercermin pada
              pohon urai.
        \item Grammar harus menghasilkan struktur yang berbeda untuk
              kedua ekspresi.
      \end{itemize}
    \end{column}
  \end{columns}
\end{frame}

% ---------------------------------------------------------------------
\begin{frame}{Kegiatan Mahasiswa di Kelas: P2-A}
  \begin{block}{Tata bahasa Subset-C}
    \footnotesize
    \begin{enumerate}
      \setlength{\itemsep}{1pt}
      \item Lengkapi \file{spec/grammar.ebnf} untuk seluruh konstruksi.
      \item Turunkan lima string: ekspresi, deklarasi, if, while, dan
            fungsi.
      \item Konversi EBNF ke BNF pada \file{spec/grammar.bnf}.
      \item Gambar pohon urai untuk dua string.
      \item Tukar berkas dengan teman untuk peer review.
    \end{enumerate}
  \end{block}

  \vspace{4pt}
  {\footnotesize\color{ubbblue!80}Periksa bahwa setiap aturan EBNF
  memiliki padanan BNF dan setiap nonterminal bantu diberi nama
  konsisten.}
\end{frame}

% ---------------------------------------------------------------------
\begin{frame}{Tugas dan Artefak}
  \begin{block}{Artefak Sub-CPMK 3.1 dan 3.4}
    Berkas grammar BNF/EBNF untuk deklarasi, ekspresi, dan pernyataan,
    disertai lima derivasi, dua pohon urai, dan catatan hasil uji.
  \end{block}

  \small
  \begin{itemize}
    \item EBNF mencakup semua konstruksi yang didukung Subset-C.
    \item Lima derivasi benar dan terdokumentasi.
    \item BNF konsisten dengan EBNF; simbol bantu dinamai konsisten.
    \item Dua pohon urai benar dan sesuai derivasi.
    \item Catatan peer review lengkap.
  \end{itemize}

  \vspace{3pt}
  {\footnotesize\color{ubbblue!75}P2-A menjadi bagian Praktikum 2
  (12\%) yang dilaporkan pada Pertemuan 9.}
\end{frame}

% ---------------------------------------------------------------------
\begin{frame}{Rangkuman dan Jembatan}
  \begin{itemize}
    \item CFG \(G=(V,T,P,S)\) membangkitkan bahasa melalui derivasi
          \(S \Rightarrow^{*} w\).
    \item BNF menulis aturan dengan alternatif; EBNF menambahkan
          pengulangan, opsional, dan pengelompokan.
    \item EBNF dikonversi ke BNF dengan aturan sistematis dan simbol
          bantu \texttt{\_tail} serta \texttt{opt\_}.
    \item Derivasi leftmost dan rightmost menghasilkan pohon urai;
          AST menyederhanakan detail sintaksis.
    \item Grammar Subset-C memisahkan deklarasi, pernyataan, dan
          tingkat prioritas ekspresi.
  \end{itemize}

  \vspace{4pt}
  \begin{exampleblock}{Pertemuan berikutnya}
    Pertemuan 7 membersihkan grammar dari ambiguitas, rekursi kiri,
    dan kebutuhan left factoring. Grammar yang bersih menjadi syarat
    parser top-down pada Pertemuan 9.
  \end{exampleblock}
\end{frame}

\end{document}
